%=====================================================================
% Sección de Belmont Muñoz Samuel: pruebas de integración
% Los valores MEDIDOS (en rojo) los llena Samuel con el osciloscopio o el
% analizador lógico. No inventar datos.
%=====================================================================
\section{Pruebas de integración}
\textbf{Autor:} Belmont Muñoz Samuel

\textbf{Requisito medido:} RF1 del Avance 1, lectura de botones con periodo de $5\text{ ms} \pm 0.5\text{ ms}$ ($200\text{ Hz}$).

\textbf{Cálculo teórico.} La lectura la hace la PIO sin intervención de la CPU, así que el periodo lo fija
el número de ciclos del programa \texttt{read\_74hc165.pio} y el divisor de reloj
(\texttt{BTN\_SM\_CLKDIV = 125}, o sea $125\text{ MHz}/125 = 1\text{ MHz}$, $1\,\mu\text{s}$ por ciclo):
\begin{align*}
	\text{carga y liberación} &= 4 \times 32 = 128 \\
	\text{primer bit} + \texttt{set x} &= 1 + 1 = 2 \\
	\text{9 bits restantes} &= 9 \times (4 \times 32 + 1 + 1) = 1170 \\
	\text{pausa} &= 1 + 2 \times (1 + 32 \times 33 + 1) = 2117 \\
	\text{total} &= 3417 \text{ ciclos} = 3.417\text{ ms}\ (292.7\text{ Hz})
\end{align*}

\textbf{Método:} la PIO ejecuta exactamente una instrucción por ciclo de su reloj, sin interrupciones ni caché
que la retrasen, así que el periodo de la trama es determinista y queda fijado por el conteo de ciclos y el divisor.
La única fuente de variación es el oscilador de cristal de 12 MHz de la Pico, del orden de partes por millón.
La medición con osciloscopio en GPIO17 (señal LD) queda como confirmación para el Avance 4.

\begin{table}[H]
	\centering
	\caption{Periodo de lectura de botones contra el Avance 1}
	\label{tab:medicion}
	\small
	\begin{tabular}{|l|c|c|c|}
		\hline
		\textbf{Dato} & \textbf{Avance 1} & \textbf{Actual} & \textbf{Ajustado} \\
		\hline
		Periodo & $5.0 \pm 0.5\text{ ms}$ & $3.417\text{ ms}$ & $5.000\text{ ms}$ \\
		\hline
		Frecuencia & $200\text{ Hz}$ & $292.7\text{ Hz}$ & $200.0\text{ Hz}$ \\
		\hline
	\end{tabular}
\end{table}

\textbf{Resultado:} con el divisor actual (\texttt{BTN\_SM\_CLKDIV = 125}) la trama dura 3.417 ms, fuera del
intervalo de $5 \pm 0.5$ ms comprometido: \textbf{no cumple}, por ser 1.583 ms más rápida. Leer más rápido no afecta
la latencia, pero el requisito pide un periodo fijo. La corrección es cambiar el divisor a 182.91:
\[
T = \frac{3417 \times 182.91}{125\text{ MHz}} = 5.000\text{ ms}
\]
El divisor fraccionario de la PIO tiene resolución de $1/256$, así que $182.91 \approx 182 + 233/256 = 182.910$,
y el error queda por debajo del 0.01\,\% (menos de $0.5\,\mu\text{s}$).
